% ID: SEL-D0030-OA-E03
% ESTADO: cotejada. Segunda pasada visual completa y borrador compilado revisados.
% Fuente: Selectividad/matematicas_aplicadas/titular_junio/Examen-Matemáticas Aplicadas a las Ciencias Sociales.pdf
% SHA-256: 5bd7838829b9f653746319907f86642017d12e66ed097db0126fb544a2d0f632
% Página: 1; opción A; ejercicio original 3.

En una determinada población residen 5000 personas en el centro y 10000 en la periferia. Se sabe que el 95\% de los residentes en el centro y que el 20\% de los que viven en la periferia opina que el Ayuntamiento debería restringir el acceso de vehículos privados al centro urbano. Se elige al azar un residente de la población.
\begin{enumerate}[label=\alph*)]
\item \textbf{(1.25 puntos)} \textquestiondown{}Cuál es la probabilidad de que esté a favor de restringir el acceso de vehículos privados al centro de la ciudad?
\item \textbf{(0.5 puntos)} \textquestiondown{}Cuál es la probabilidad de que resida en el centro y esté a favor de la restricción de acceso?
\item \textbf{(0.75 puntos)} Si la persona elegida opina que se debería restringir el acceso, \textquestiondown{}cuál es la probabilidad de que resida en el centro de la ciudad?
\end{enumerate}
